\subsection{HDT as a Promising Solution}

\begin{frame}{What is studied?}
    \begin{columns}
        \column{0.6\textwidth}
            \begin{itemize}
                \only<1-2>{
                \item A Hybrid distribution transformer (HDT) can be used to improve the power quality of a distribution network and to extend the lifespan of the distribution transformer (DT).
                \item The way to make control to the HDT is separated.
                \begin{itemize}
                    \item One LQR control for the series converter.
                    \item One multi-resonant state-feedback control for the parallel converter.
                \end{itemize}
                }
                \only<2>{
                \begin{exampleblock}{Observations}
                    \begin{itemize}
                        \item This can be unified in only one controller (which could be called "Internal Control Loop").
                    \end{itemize}
                \end{exampleblock}
                }
                \only<3-4>{
                \item The HDT issue of circulating active power flow (CAPF) can be eliminated if a DC microgrid (DC-MG) is connected to the DC-Link of both voltage and current converters.
                \item This enables the HDT to act as an interlinking converter (IC) between these two MGs.
                }
                \only<4>{
                \begin{exampleblock}{Observations}
                    \begin{itemize}
                        \item This allows the IC to be controlled by a hierarchical control for MGs (which will have the internal control nested to it).
                    \end{itemize}
                \end{exampleblock}
                }
            \end{itemize}
        \column{0.4\textwidth}
            \begin{figure}
                \centering
                \resizebox{!}{5.3cm}{

\tikzset{every picture/.style={line width=0.75pt}} %set default line width to 0.75pt        

\begin{tikzpicture}[x=0.75pt,y=0.75pt,yscale=-1,xscale=1]
%uncomment if require: \path (0,281); %set diagram left start at 0, and has height of 281

%Rounded Rect [id:dp1498416922615955] 
\draw  [fill={rgb, 255:red, 255; green, 0; blue, 0 }  ,fill opacity=0.2 ] (48,135.75) .. controls (48,131.47) and (51.47,128) .. (55.75,128) -- (184.25,128) .. controls (188.53,128) and (192,131.47) .. (192,135.75) -- (192,248.25) .. controls (192,252.53) and (188.53,256) .. (184.25,256) -- (55.75,256) .. controls (51.47,256) and (48,252.53) .. (48,248.25) -- cycle ;
%Shape: Triangle [id:dp46443946003277636] 
\draw  [fill={rgb, 255:red, 255; green, 255; blue, 255 }  ,fill opacity=1 ] (119.41,63.91) -- (95.81,87.95) -- (95.81,39.87) -- cycle ;
%Shape: Polygon [id:ds8300633911306601] 
\draw  [fill={rgb, 255:red, 255; green, 255; blue, 255 }  ,fill opacity=1 ][general shadow={fill=black,shadow xshift=0.75pt,shadow yshift=-0.75pt}] (143,39.87) -- (143,87.95) -- (97.29,87.95) -- (120.88,63.91) -- (97.29,39.87) -- cycle ;
%Straight Lines [id:da7060653075357857] 
\draw    (120.88,63.91) -- (143,63.91) ;

%Straight Lines [id:da6363435803598962] 
\draw    (64,64) -- (96,64) ;
%Straight Lines [id:da5267996923036957] 
\draw    (120,88) -- (120,108) ;
%Straight Lines [id:da44726021460361265] 
\draw    (120,8) -- (120,40) ;
%Straight Lines [id:da668258115881776] 
\draw    (95.81,87.95) -- (56,128) ;
%Straight Lines [id:da9214301838509702] 
\draw    (144,88) -- (184,128) ;
%Shape: Inductor (Air Core) [id:dp83632344975283] 
\draw   (88,152) -- (83.68,152) .. controls (83.68,153.77) and (82.82,155.2) .. (81.76,155.2) .. controls (80.7,155.2) and (79.84,153.77) .. (79.84,152) .. controls (79.84,153.77) and (78.98,155.2) .. (77.92,155.2) .. controls (76.86,155.2) and (76,153.77) .. (76,152) .. controls (76,153.77) and (75.14,155.2) .. (74.08,155.2) .. controls (73.02,155.2) and (72.16,153.77) .. (72.16,152) .. controls (72.16,153.77) and (71.3,155.2) .. (70.24,155.2) .. controls (69.18,155.2) and (68.32,153.77) .. (68.32,152) -- (64,152) ;
%Shape: Inductor (Air Core) [id:dp7999508751202471] 
\draw   (64,168) -- (68.32,168) .. controls (68.32,166.23) and (69.18,164.8) .. (70.24,164.8) .. controls (71.3,164.8) and (72.16,166.23) .. (72.16,168) .. controls (72.16,166.23) and (73.02,164.8) .. (74.08,164.8) .. controls (75.14,164.8) and (76,166.23) .. (76,168) .. controls (76,166.23) and (76.86,164.8) .. (77.92,164.8) .. controls (78.98,164.8) and (79.84,166.23) .. (79.84,168) .. controls (79.84,166.23) and (80.7,164.8) .. (81.76,164.8) .. controls (82.82,164.8) and (83.68,166.23) .. (83.68,168) -- (88,168) ;
%Straight Lines [id:da8313103681431602] 
\draw    (85,158.4) -- (67,158.4) ;
%Straight Lines [id:da41157882719633565] 
\draw    (85,161.6) -- (67,161.6) ;

%Shape: Ellipse [id:dp9743148466945699] 
\draw  [fill={rgb, 255:red, 255; green, 255; blue, 255 }  ,fill opacity=1 ] (108.32,152) .. controls (108.32,147.58) and (111.81,144) .. (116.11,144) .. controls (120.41,144) and (123.89,147.58) .. (123.89,152) .. controls (123.89,156.42) and (120.41,160) .. (116.11,160) .. controls (111.81,160) and (108.32,156.42) .. (108.32,152) -- cycle ;
%Straight Lines [id:da3385079185306148] 
\draw    (88,152) -- (108.32,152) ;
%Shape: Ellipse [id:dp4056578022596924] 
\draw  [fill={rgb, 255:red, 255; green, 255; blue, 255 }  ,fill opacity=1 ] (120,152) .. controls (120,147.58) and (123.48,144) .. (127.78,144) .. controls (132.08,144) and (135.57,147.58) .. (135.57,152) .. controls (135.57,156.42) and (132.08,160) .. (127.78,160) .. controls (123.48,160) and (120,156.42) .. (120,152) -- cycle ;
%Shape: Rectangle [id:dp5537843893958982] 
\draw  [fill={rgb, 255:red, 255; green, 255; blue, 255 }  ,fill opacity=1 ][general shadow={fill=black,shadow xshift=0.75pt,shadow yshift=-0.75pt}] (88.23,200.51) -- (111.09,200.51) -- (111.09,223.76) -- (88.23,223.76) -- cycle ;
%Straight Lines [id:da45439810694201177] 
\draw    (111.09,200.51) -- (88.23,223.76) ;

%Shape: Rectangle [id:dp3659605277693705] 
\draw  [fill={rgb, 255:red, 255; green, 255; blue, 255 }  ,fill opacity=1 ][general shadow={fill=black,shadow xshift=0.75pt,shadow yshift=-0.75pt}] (136.55,200.51) -- (159.66,200.51) -- (159.66,223.76) -- (136.55,223.76) -- cycle ;
%Straight Lines [id:da9862335388880625] 
\draw    (159.66,200.51) -- (136.55,223.76) ;

%Straight Lines [id:da8530054247718197] 
\draw    (76,172) -- (76,212) -- (87.68,212) ;
%Shape: Capacitor [id:dp052526262560992576] 
\draw   (127.57,204) -- (127.57,211.2) (131.46,212.8) -- (123.68,212.8) (131.46,211.2) -- (123.68,211.2) (127.57,212.8) -- (127.57,220) ;
%Straight Lines [id:da06685240362324807] 
\draw    (112,204) -- (128,204) ;
\draw [shift={(128,204)}, rotate = 0] [color={rgb, 255:red, 0; green, 0; blue, 0 }  ][fill={rgb, 255:red, 0; green, 0; blue, 0 }  ][line width=0.75]      (0, 0) circle [x radius= 1.34, y radius= 1.34]   ;
%Straight Lines [id:da08687878975255625] 
\draw    (128,204) -- (136,204) ;
%Straight Lines [id:da05664366026061396] 
\draw    (112,220) -- (127.57,220) ;
\draw [shift={(127.57,220)}, rotate = 0] [color={rgb, 255:red, 0; green, 0; blue, 0 }  ][fill={rgb, 255:red, 0; green, 0; blue, 0 }  ][line width=0.75]      (0, 0) circle [x radius= 1.34, y radius= 1.34]   ;
%Straight Lines [id:da18202214085786328] 
\draw    (127.57,220) -- (136,220) ;
%Straight Lines [id:da8257196245918859] 
\draw    (172,152) -- (172,212) -- (160,212) ;
\draw [shift={(172,152)}, rotate = 90] [color={rgb, 255:red, 0; green, 0; blue, 0 }  ][fill={rgb, 255:red, 0; green, 0; blue, 0 }  ][line width=0.75]      (0, 0) circle [x radius= 1.34, y radius= 1.34]   ;
%Straight Lines [id:da15625609577334676] 
\draw    (135.57,152) -- (176.43,152) ;
%Straight Lines [id:da7061059641709855] 
\draw    (50.35,152) -- (64,152) ;
\draw [shift={(48,152)}, rotate = 0] [color={rgb, 255:red, 0; green, 0; blue, 0 }  ][line width=0.75]      (0, 0) circle [x radius= 3.35, y radius= 3.35]   ;
%Straight Lines [id:da7563615687380554] 
\draw    (176.43,152) -- (189.65,152) ;
\draw [shift={(192,152)}, rotate = 0] [color={rgb, 255:red, 0; green, 0; blue, 0 }  ][line width=0.75]      (0, 0) circle [x radius= 3.35, y radius= 3.35]   ;
%Straight Lines [id:da9257628912880724] 
\draw    (120,204) -- (120,253.65) ;
\draw [shift={(120,256)}, rotate = 90] [color={rgb, 255:red, 0; green, 0; blue, 0 }  ][line width=0.75]      (0, 0) circle [x radius= 3.35, y radius= 3.35]   ;
%Shape: Simple Switch [id:dp3562117821380484] 
\draw   (68,148) -- (71.2,148) (80.8,148) -- (84,148) (72.48,147.75) -- (80.16,143.98) (79.52,148) .. controls (79.52,147.58) and (79.81,147.25) .. (80.16,147.25) .. controls (80.51,147.25) and (80.8,147.58) .. (80.8,148) .. controls (80.8,148.42) and (80.51,148.75) .. (80.16,148.75) .. controls (79.81,148.75) and (79.52,148.42) .. (79.52,148) -- cycle (71.2,148) .. controls (71.2,147.58) and (71.49,147.25) .. (71.84,147.25) .. controls (72.19,147.25) and (72.48,147.58) .. (72.48,148) .. controls (72.48,148.42) and (72.19,148.75) .. (71.84,148.75) .. controls (71.49,148.75) and (71.2,148.42) .. (71.2,148) -- cycle ;
%Straight Lines [id:da8570241382152515] 
\draw    (68,148) -- (64,148) -- (64,152) ;
\draw [shift={(64,152)}, rotate = 90] [color={rgb, 255:red, 0; green, 0; blue, 0 }  ][fill={rgb, 255:red, 0; green, 0; blue, 0 }  ][line width=0.75]      (0, 0) circle [x radius= 1.34, y radius= 1.34]   ;
%Straight Lines [id:da8968836011702404] 
\draw    (84,148) -- (88,148) -- (88,152) ;
\draw [shift={(88,152)}, rotate = 90] [color={rgb, 255:red, 0; green, 0; blue, 0 }  ][fill={rgb, 255:red, 0; green, 0; blue, 0 }  ][line width=0.75]      (0, 0) circle [x radius= 1.34, y radius= 1.34]   ;

% Text Node
\draw (119.75,1.5) node  [font=\small] [align=left] {\textbf{LV-AC}};
% Text Node
\draw (120,117.5) node  [font=\small] [align=left] {\textbf{LV-DC}};
% Text Node
\draw (41,54.5) node [anchor=west] [inner sep=0.75pt]  [font=\small] [align=left] {\textbf{MV-DC}};
% Text Node
\draw (134.5,47.5) node  [font=\scriptsize] [align=left] {AC};
% Text Node
\draw (105.5,63.5) node  [font=\scriptsize] [align=left] {AC};
% Text Node
\draw (134.5,80.5) node  [font=\scriptsize] [align=left] {DC};
% Text Node
\draw (116.05,152) node  [font=\tiny]  {$\Delta $};
% Text Node
\draw (127.7,152) node  [font=\tiny]  {$Y$};
% Text Node
\draw (95.14,205.5) node  [font=\tiny] [align=left] {AC};
% Text Node
\draw (104.35,218.5) node  [font=\tiny] [align=left] {DC};
% Text Node
\draw (152.86,218.5) node  [font=\tiny] [align=left] {AC};
% Text Node
\draw (143.54,205.5) node  [font=\tiny] [align=left] {DC};
% Text Node
\draw (82.5,237.5) node  [font=\tiny] [align=left] {Voltage Controlled\\Power Converter};
% Text Node
\draw (158,237.5) node  [font=\tiny] [align=left] {Current Controller\\Power Converter};
% Text Node
\draw (122.5,194.5) node  [font=\tiny] [align=left] {DC-Link};
% Text Node
\draw (48.62,142.5) node [anchor=west] [inner sep=0.75pt]  [font=\tiny] [align=left] {\textbf{MV-AC}};
% Text Node
\draw (191.35,145.5) node [anchor=east] [inner sep=0.75pt]  [font=\tiny] [align=left] {\textbf{LV-AC}};
% Text Node
\draw (120,265.5) node  [font=\tiny] [align=left] {\textbf{LV-DC}};


\end{tikzpicture}}
                \caption{Hybrid distribution Transformer.}
                \label{fig:HDT}
            \end{figure}
    \end{columns}
\end{frame}

\begin{frame}{Hierarchical Control of Microgrids}
    \begin{columns}
        \column{0.5\textwidth}
            \only<1>{
                \begin{block}{Primary Control}
                    \begin{itemize}
                        \item Stabilize the amplitude voltage and frequency of the microgrid.
                        \item Provide capacity for the distributed energy resource (DER) to connect and disconnect from the grid.
                        \item Protect the converter.
                    \end{itemize}
                \end{block}
            }
            \only<2>{
                \begin{block}{Secondary Control}
                    \begin{itemize}
                        \item Gives the nominal operating point for the primary control.
                        \item Provides the ability for the DER to contribute active power to the main grid. 
                    \end{itemize}
                \end{block}
            }
            \only<3>{
                \begin{block}{Tertiary Control}
                    \begin{itemize}
                        \item Coordinates power dispatch to the main grid.
                    \end{itemize}
                \end{block}
            }
        \column{0.5\textwidth}
            \begin{figure}
                \centering
                \resizebox{!}{4.5cm}{
                \only<1>{\input{Images/HierarchicalMG-PC}}
                \only<2>{\input{Images/HierarchicalMG-SC}}
                \only<3>{

\tikzset{every picture/.style={line width=0.75pt}} %set default line width to 0.75pt        

\begin{tikzpicture}[x=0.75pt,y=0.75pt,yscale=-1,xscale=1]
%uncomment if require: \path (0,302); %set diagram left start at 0, and has height of 302

%Shape: Rectangle [id:dp9309347763461011] 
\draw  [fill={rgb, 255:red, 255; green, 255; blue, 255 } ,fill opacity=1 ][general shadow={fill=black,shadow xshift=1.5pt,shadow yshift=-1.5pt}] (285,128) -- (407,128) -- (407,156) -- (285,156) -- cycle ;
%Shape: Rectangle [id:dp9321591279013901] 
\draw  [fill={rgb, 255:red, 255; green, 255; blue, 255 }  ,fill opacity=1 ][general shadow={fill=black,shadow xshift=1.5pt,shadow yshift=-1.5pt}] (263,80) -- (391,80) -- (391,108) -- (263,108) -- cycle ;
%Shape: Rectangle [id:dp15063633321617553] 
\draw  [fill={rgb, 255:red, 253; green, 192; blue, 0 }  ,fill opacity=1 ][general shadow={fill=black,shadow xshift=1.5pt,shadow yshift=-1.5pt}] (208,8) -- (344,8) -- (344,40) -- (208,40) -- cycle ;
%Shape: Rectangle [id:dp4651413853898998] 
\draw  [fill={rgb, 255:red, 255; green, 255; blue, 255 }  ,fill opacity=1 ][general shadow={fill=black,shadow xshift=1.5pt,shadow yshift=-1.5pt}] (232,192) -- (416,192) -- (416,224) -- (232,224) -- cycle ;
%Shape: Rectangle [id:dp11630061840392214] 
\draw  [fill={rgb, 255:red, 255; green, 255; blue, 255 }  ,fill opacity=1 ][general shadow={fill=black,shadow xshift=1.5pt,shadow yshift=-1.5pt}] (208,256) -- (432,256) -- (432,288) -- (208,288) -- cycle ;
%Straight Lines [id:da8520220149226383] 
\draw [color={rgb, 255:red, 65; green, 117; blue, 5 }  ,draw opacity=1 ]   (224,256) -- (224,43) ;
\draw [shift={(224,40)}, rotate = 90] [fill={rgb, 255:red, 65; green, 117; blue, 5 }  ,fill opacity=1 ][line width=0.08]  [draw opacity=0] (8.93,-4.29) -- (0,0) -- (8.93,4.29) -- cycle    ;
%Straight Lines [id:da018327241437618103] 
\draw [color={rgb, 255:red, 65; green, 117; blue, 5 }  ,draw opacity=1 ]   (240,192) -- (240,43) ;
\draw [shift={(240,40)}, rotate = 90] [fill={rgb, 255:red, 65; green, 117; blue, 5 }  ,fill opacity=1 ][line width=0.08]  [draw opacity=0] (8.93,-4.29) -- (0,0) -- (8.93,4.29) -- cycle    ;
%Straight Lines [id:da14281077756300853] 
\draw [color={rgb, 255:red, 65; green, 117; blue, 5 }  ,draw opacity=1 ]   (272,192) -- (272,115) ;
\draw [shift={(272,112)}, rotate = 90] [fill={rgb, 255:red, 65; green, 117; blue, 5 }  ,fill opacity=1 ][line width=0.08]  [draw opacity=0] (8.93,-4.29) -- (0,0) -- (8.93,4.29) -- cycle    ;
%Straight Lines [id:da67579633614294] 
\draw [color={rgb, 255:red, 65; green, 117; blue, 5 }  ,draw opacity=1 ]   (311,192) -- (311,171.09) -- (311,160) ;
\draw [shift={(311,157)}, rotate = 90] [fill={rgb, 255:red, 65; green, 117; blue, 5 }  ,fill opacity=1 ][line width=0.08]  [draw opacity=0] (8.93,-4.29) -- (0,0) -- (8.93,4.29) -- cycle    ;
%Straight Lines [id:da21350666841290278] 
\draw [color={rgb, 255:red, 65; green, 117; blue, 5 }  ,draw opacity=1 ]   (327,109) -- (327,125) ;
\draw [shift={(327,128)}, rotate = 270] [fill={rgb, 255:red, 65; green, 117; blue, 5 }  ,fill opacity=1 ][line width=0.08]  [draw opacity=0] (8.93,-4.29) -- (0,0) -- (8.93,4.29) -- cycle    ;
%Straight Lines [id:da7751555879812528] 
\draw [color={rgb, 255:red, 65; green, 117; blue, 5 }  ,draw opacity=1 ]   (319,40) -- (319,77) ;
\draw [shift={(319,80)}, rotate = 270] [fill={rgb, 255:red, 65; green, 117; blue, 5 }  ,fill opacity=1 ][line width=0.08]  [draw opacity=0] (8.93,-4.29) -- (0,0) -- (8.93,4.29) -- cycle    ;
%Straight Lines [id:da8361685456904682] 
\draw [color={rgb, 255:red, 65; green, 117; blue, 5 }  ,draw opacity=1 ]   (383,160) -- (383,189) ;
\draw [shift={(383,192)}, rotate = 270] [fill={rgb, 255:red, 65; green, 117; blue, 5 }  ,fill opacity=1 ][line width=0.08]  [draw opacity=0] (8.93,-4.29) -- (0,0) -- (8.93,4.29) -- cycle    ;
%Straight Lines [id:da7956201861715562] 
\draw    (328,226) -- (328,256) ;

% Text Node
\draw (346,142) node  [font=\footnotesize] [align=left] {\textbf{Primary Control}};
% Text Node
\draw (326.5,94) node  [font=\footnotesize] [align=left] {\textbf{Secondary Control}};
% Text Node
\draw (276,24) node  [font=\footnotesize] [align=left] {\textbf{Tertiary Control}};
% Text Node
\draw (324,208) node   [align=left] {\textbf{Local Grid o Microgrid}};
% Text Node
\draw (320,272) node   [align=left] {\textbf{Grid}};
% Text Node
\draw  [draw opacity=0][fill={rgb, 255:red, 255; green, 255; blue, 255 }  ,fill opacity=1 ]  (215.4,161.6) .. controls (215.4,161.05) and (215.85,160.6) .. (216.4,160.6) -- (283.4,160.6) .. controls (283.95,160.6) and (284.4,161.05) .. (284.4,161.6) -- (284.4,180.6) .. controls (284.4,181.15) and (283.95,181.6) .. (283.4,181.6) -- (216.4,181.6) .. controls (215.85,181.6) and (215.4,181.15) .. (215.4,180.6) -- cycle  ;
\draw (218.4,164.6) node [anchor=north west][inner sep=0.75pt]  [font=\footnotesize] [align=left] {Measurement};
% Text Node
\draw (320,50) node [anchor=north west][inner sep=0.75pt]   [align=left] {Control Action};


\end{tikzpicture}}
                }
                \caption{Microgrid hierarchical control.}
                \label{fig:MGHierarchy}
        \end{figure}
    \end{columns}
\end{frame}

\begin{frame}{Other Ways to Make Control of the Internal Loop of IC}
    \begin{itemize}
        \item Model predictive control (MPC)
        \begin{itemize}
            \item It is a control strategy that uses a model of the system to predict its future behavior and optimize the control inputs accordingly.
            \item One of the key advantages of MPC is its ability to impose restrictions on either the states and/or the control inputs.
            \item This makes it useful if the objective is to protect the IC.
        \end{itemize}
        \item Reinforcement learning (RL)
        \begin{itemize}
            \item It is a type of machine learning that focuses on how agents should take actions in an environment to maximize a cumulative reward.
            \item RL has been successfully applied to various control problems, including power systems.
            \item The main advantage of RL is its ability to learn from experience and adapt to changing conditions.
        \end{itemize}
    \end{itemize}
\end{frame}

\begin{frame}{HDTs as a Promising Solution for AC and DC Microgrids Integration}
    \begin{figure}
        \centering
        \resizebox{!}{5.3cm}{

\tikzset{every picture/.style={line width=0.75pt}} %set default line width to 0.75pt        

\begin{tikzpicture}[x=0.75pt,y=0.75pt,yscale=-1,xscale=1]
%uncomment if require: \path (0,315); %set diagram left start at 0, and has height of 315

%Shape: Rectangle [id:dp36299639453955] 
\draw  [color={rgb, 255:red, 255; green, 0; blue, 0 }  ,draw opacity=1 ][fill={rgb, 255:red, 255; green, 0; blue, 0 }  ,fill opacity=0.25 ][dash pattern={on 4.5pt off 4.5pt}] (88,133) .. controls (88,130.24) and (90.24,128) .. (93,128) -- (187,128) .. controls (189.76,128) and (192,130.24) .. (192,133) -- (192,171) .. controls (192,173.76) and (189.76,176) .. (187,176) -- (93,176) .. controls (90.24,176) and (88,173.76) .. (88,171) -- cycle ;
%Shape: Rectangle [id:dp09007032268016424] 
\draw  [color={rgb, 255:red, 126; green, 211; blue, 33 }  ,draw opacity=1 ][fill={rgb, 255:red, 126; green, 211; blue, 33 }  ,fill opacity=0.25 ] (200,21) .. controls (200,18.24) and (202.24,16) .. (205,16) -- (243,16) .. controls (245.76,16) and (248,18.24) .. (248,21) -- (248,51) .. controls (248,53.76) and (245.76,56) .. (243,56) -- (205,56) .. controls (202.24,56) and (200,53.76) .. (200,51) -- cycle ;
%Shape: Rectangle [id:dp7538822078187675] 
\draw  [color={rgb, 255:red, 80; green, 227; blue, 194 }  ,draw opacity=1 ][fill={rgb, 255:red, 80; green, 227; blue, 194 }  ,fill opacity=0.25 ] (144,21) .. controls (144,18.24) and (146.24,16) .. (149,16) -- (187,16) .. controls (189.76,16) and (192,18.24) .. (192,21) -- (192,51) .. controls (192,53.76) and (189.76,56) .. (187,56) -- (149,56) .. controls (146.24,56) and (144,53.76) .. (144,51) -- cycle ;
%Shape: Rectangle [id:dp38254314765596953] 
\draw   (152,32) -- (184,32) -- (184,48) -- (152,48) -- cycle ;
%Shape: Rectangle [id:dp40775062639027393] 
\draw  [fill={rgb, 255:red, 0; green, 0; blue, 0 }  ,fill opacity=1 ] (154,30) -- (158,30) -- (158,32) -- (154,32) -- cycle ;
%Shape: Rectangle [id:dp8484922427356605] 
\draw  [fill={rgb, 255:red, 0; green, 0; blue, 0 }  ,fill opacity=1 ] (178,30) -- (182,30) -- (182,32) -- (178,32) -- cycle ;
%Straight Lines [id:da9811935771636653] 
\draw    (154,36) -- (158,36) ;
%Straight Lines [id:da23592258919271258] 
\draw    (178,36) -- (182,36) ;
%Straight Lines [id:da17310904483441059] 
\draw    (180,34) -- (180,38) ;


%Shape: Rectangle [id:dp608370505344022] 
\draw  [color={rgb, 255:red, 245; green, 166; blue, 35 }  ,draw opacity=1 ][fill={rgb, 255:red, 245; green, 166; blue, 35 }  ,fill opacity=0.25 ] (88,21) .. controls (88,18.24) and (90.24,16) .. (93,16) -- (131,16) .. controls (133.76,16) and (136,18.24) .. (136,21) -- (136,51) .. controls (136,53.76) and (133.76,56) .. (131,56) -- (93,56) .. controls (90.24,56) and (88,53.76) .. (88,51) -- cycle ;
%Shape: Parallelogram [id:dp5542004190325218] 
\draw  [fill={rgb, 255:red, 0; green, 0; blue, 255 }  ,fill opacity=0.25 ] (100.1,24) -- (132,24) -- (123.9,40) -- (92,40) -- cycle ;
%Straight Lines [id:da5044587363845592] 
\draw    (96,40) -- (100,32) ;
%Straight Lines [id:da007739550083210478] 
\draw    (100,40) -- (104,32) ;
%Straight Lines [id:da8618146610691602] 
\draw    (104,40) -- (108,32) ;
%Straight Lines [id:da3063231199330758] 
\draw    (108,40) -- (112,32) ;
%Straight Lines [id:da3873698645587713] 
\draw    (112,40) -- (116,32) ;
%Straight Lines [id:da07382503135839613] 
\draw    (116,40) -- (120,32) ;
%Straight Lines [id:da6307749528038389] 
\draw    (120,40) -- (124,32) ;
%Straight Lines [id:da7078011067031311] 
\draw    (94,36) -- (126,36) ;
%Straight Lines [id:da36003974470731626] 
\draw    (100,32) -- (104,24) ;
%Straight Lines [id:da6574946945911684] 
\draw    (104,32) -- (108,24) ;
%Straight Lines [id:da6630449731731625] 
\draw    (108,32) -- (112,24) ;
%Straight Lines [id:da7804944245776355] 
\draw    (112,32) -- (116,24) ;
%Straight Lines [id:da8297470482526796] 
\draw    (116,32) -- (120,24) ;
%Straight Lines [id:da7048240489300945] 
\draw    (120,32) -- (124,24) ;
%Straight Lines [id:da9645536648637945] 
\draw    (124,32) -- (128,24) ;
%Straight Lines [id:da4231587866518627] 
\draw    (98,28) -- (130,28) ;
%Shape: Rectangle [id:dp49159466073928093] 
\draw  [fill={rgb, 255:red, 0; green, 0; blue, 0 }  ,fill opacity=1 ] (100,46) -- (124,46) -- (124,48) -- (100,48) -- cycle ;
%Shape: Rectangle [id:dp9010988432626634] 
\draw  [fill={rgb, 255:red, 0; green, 0; blue, 0 }  ,fill opacity=1 ] (110,40) -- (114,40) -- (114,46) -- (110,46) -- cycle ;
%Straight Lines [id:da18563016734416338] 
\draw    (96,32) -- (128,32) ;

%Shape: Rectangle [id:dp5266278023932072] 
\draw  [fill={rgb, 255:red, 255; green, 255; blue, 255 }  ,fill opacity=1 ][general shadow={fill=black,shadow xshift=0.75pt,shadow yshift=-0.75pt}] (95.85,64.53) -- (128,64.53) -- (128,96) -- (95.85,96) -- cycle ;
%Straight Lines [id:da09304750451052368] 
\draw    (128,64) -- (95.85,97.3) ;

%Straight Lines [id:da17655468696959464] 
\draw [line width=2.25]    (72,112) -- (256,112) ;
%Shape: Ellipse [id:dp7656190687221771] 
\draw   (41,152) .. controls (41,145.37) and (46.22,140) .. (52.65,140) .. controls (59.09,140) and (64.3,145.37) .. (64.3,152) .. controls (64.3,158.63) and (59.09,164) .. (52.65,164) .. controls (46.22,164) and (41,158.63) .. (41,152) -- cycle ;
%Curve Lines [id:da6261869691155981] 
\draw    (45.62,151.26) .. controls (52.47,132.78) and (50.8,168.78) .. (59.31,151.64) ;

%Straight Lines [id:da8271638976646392] 
\draw    (112,56) -- (112,64) ;
%Shape: Rectangle [id:dp4100675503446629] 
\draw  [fill={rgb, 255:red, 255; green, 255; blue, 255 }  ,fill opacity=1 ][general shadow={fill=black,shadow xshift=0.75pt,shadow yshift=-0.75pt}] (152,64.53) -- (184.15,64.53) -- (184.15,96) -- (152,96) -- cycle ;
%Straight Lines [id:da6292272341172389] 
\draw    (184.15,64) -- (152,97.3) ;

%Straight Lines [id:da12174681331026616] 
\draw    (168,56) -- (168,64) ;
%Straight Lines [id:da4329994692833059] 
\draw    (112,96) -- (112,112) ;
\draw [shift={(112,112)}, rotate = 90] [color={rgb, 255:red, 0; green, 0; blue, 0 }  ][fill={rgb, 255:red, 0; green, 0; blue, 0 }  ][line width=0.75]      (0, 0) circle [x radius= 2.01, y radius= 2.01]   ;
%Straight Lines [id:da5644216395705126] 
\draw    (168,96) -- (168,112) ;
\draw [shift={(168,112)}, rotate = 90] [color={rgb, 255:red, 0; green, 0; blue, 0 }  ][fill={rgb, 255:red, 0; green, 0; blue, 0 }  ][line width=0.75]      (0, 0) circle [x radius= 2.01, y radius= 2.01]   ;
%Shape: Path Data [id:dp9539730007897249] 
\draw  [fill={rgb, 255:red, 245; green, 166; blue, 35 }  ,fill opacity=0.25 ] (240,48) -- (228,48) -- (228,40) -- (220,40) -- (220,48) -- (208,48) -- (208,32) -- (224,24) -- (240,32) -- (240,48) -- cycle ;
%Straight Lines [id:da8603794313537159] 
\draw    (204,32) -- (224,22) ;
%Straight Lines [id:da8962214948125011] 
\draw    (224,22) -- (244,32) ;
%Shape: Rectangle [id:dp453727161992874] 
\draw  [fill={rgb, 255:red, 0; green, 0; blue, 0 }  ,fill opacity=1 ] (236,24) -- (238,24) -- (238,30) -- (236,30) -- cycle ;

%Straight Lines [id:da5355085247991662] 
\draw    (224,56) -- (224,112) ;
\draw [shift={(224,112)}, rotate = 90] [color={rgb, 255:red, 0; green, 0; blue, 0 }  ][fill={rgb, 255:red, 0; green, 0; blue, 0 }  ][line width=0.75]      (0, 0) circle [x radius= 2.01, y radius= 2.01]   ;
%Shape: Rectangle [id:dp38730832406139215] 
\draw  [color={rgb, 255:red, 126; green, 211; blue, 33 }  ,draw opacity=1 ][fill={rgb, 255:red, 126; green, 211; blue, 33 }  ,fill opacity=0.25 ] (200,253) .. controls (200,250.24) and (202.24,248) .. (205,248) -- (243,248) .. controls (245.76,248) and (248,250.24) .. (248,253) -- (248,283) .. controls (248,285.76) and (245.76,288) .. (243,288) -- (205,288) .. controls (202.24,288) and (200,285.76) .. (200,283) -- cycle ;
%Shape: Circle [id:dp8211119558841091] 
\draw   (206,262) .. controls (206,258.69) and (208.69,256) .. (212,256) .. controls (215.31,256) and (218,258.69) .. (218,262) .. controls (218,265.31) and (215.31,268) .. (212,268) .. controls (208.69,268) and (206,265.31) .. (206,262) -- cycle ;
%Shape: Circle [id:dp9453082981808478] 
\draw   (208,261) .. controls (208,260.45) and (208.45,260) .. (209,260) .. controls (209.55,260) and (210,260.45) .. (210,261) .. controls (210,261.55) and (209.55,262) .. (209,262) .. controls (208.45,262) and (208,261.55) .. (208,261) -- cycle ;
%Shape: Circle [id:dp4615825814572241] 
\draw   (214,261) .. controls (214,260.45) and (214.45,260) .. (215,260) .. controls (215.55,260) and (216,260.45) .. (216,261) .. controls (216,261.55) and (215.55,262) .. (215,262) .. controls (214.45,262) and (214,261.55) .. (214,261) -- cycle ;
%Shape: Circle [id:dp07473702214392075] 
\draw   (211,265) .. controls (211,264.45) and (211.45,264) .. (212,264) .. controls (212.55,264) and (213,264.45) .. (213,265) .. controls (213,265.55) and (212.55,266) .. (212,266) .. controls (211.45,266) and (211,265.55) .. (211,265) -- cycle ;

%Shape: Rectangle [id:dp13930109452257633] 
\draw   (224,256) -- (240,256) -- (240,280) -- (224,280) -- cycle ;
%Curve Lines [id:da2056385298136112] 
\draw    (212,268) .. controls (209.25,278.92) and (220.45,285.69) .. (224,264) ;
%Shape: Rectangle [id:dp8504838713938807] 
\draw  [fill={rgb, 255:red, 155; green, 155; blue, 155 }  ,fill opacity=1 ] (210,268) -- (214,268) -- (214,274) -- (210,274) -- cycle ;
%Shape: Rectangle [id:dp5317465538337975] 
\draw  [fill={rgb, 255:red, 255; green, 0; blue, 0 }  ,fill opacity=0.25 ] (226,258) -- (238,258) -- (238,264) -- (226,264) -- cycle ;

%Shape: Rectangle [id:dp3203086254752263] 
\draw  [color={rgb, 255:red, 80; green, 227; blue, 194 }  ,draw opacity=1 ][fill={rgb, 255:red, 80; green, 227; blue, 194 }  ,fill opacity=0.25 ] (144,253) .. controls (144,250.24) and (146.24,248) .. (149,248) -- (187,248) .. controls (189.76,248) and (192,250.24) .. (192,253) -- (192,283) .. controls (192,285.76) and (189.76,288) .. (187,288) -- (149,288) .. controls (146.24,288) and (144,285.76) .. (144,283) -- cycle ;
%Shape: Rectangle [id:dp2653700015147502] 
\draw   (152,264) -- (184,264) -- (184,280) -- (152,280) -- cycle ;
%Shape: Rectangle [id:dp8930820672294637] 
\draw  [fill={rgb, 255:red, 0; green, 0; blue, 0 }  ,fill opacity=1 ] (154,262) -- (158,262) -- (158,264) -- (154,264) -- cycle ;
%Shape: Rectangle [id:dp5896749728034121] 
\draw  [fill={rgb, 255:red, 0; green, 0; blue, 0 }  ,fill opacity=1 ] (178,262) -- (182,262) -- (182,264) -- (178,264) -- cycle ;
%Straight Lines [id:da7887351195081793] 
\draw    (154,268) -- (158,268) ;
%Straight Lines [id:da5561146518517932] 
\draw    (178,268) -- (182,268) ;
%Straight Lines [id:da7643548865379168] 
\draw    (180,266) -- (180,270) ;


%Shape: Rectangle [id:dp3575939081610673] 
\draw  [color={rgb, 255:red, 245; green, 166; blue, 35 }  ,draw opacity=1 ][fill={rgb, 255:red, 245; green, 166; blue, 35 }  ,fill opacity=0.25 ] (88,253) .. controls (88,250.24) and (90.24,248) .. (93,248) -- (131,248) .. controls (133.76,248) and (136,250.24) .. (136,253) -- (136,283) .. controls (136,285.76) and (133.76,288) .. (131,288) -- (93,288) .. controls (90.24,288) and (88,285.76) .. (88,283) -- cycle ;
%Shape: Parallelogram [id:dp8288201139793847] 
\draw  [fill={rgb, 255:red, 0; green, 0; blue, 255 }  ,fill opacity=0.25 ] (100.1,256) -- (132,256) -- (123.9,272) -- (92,272) -- cycle ;
%Straight Lines [id:da23574909008404643] 
\draw    (96,272) -- (100,264) ;
%Straight Lines [id:da4671840220374326] 
\draw    (100,272) -- (104,264) ;
%Straight Lines [id:da5542049138514937] 
\draw    (104,272) -- (108,264) ;
%Straight Lines [id:da45593053792004445] 
\draw    (108,272) -- (112,264) ;
%Straight Lines [id:da5890467686626477] 
\draw    (112,272) -- (116,264) ;
%Straight Lines [id:da7639730134684768] 
\draw    (116,272) -- (120,264) ;
%Straight Lines [id:da6457526609453472] 
\draw    (120,272) -- (124,264) ;
%Straight Lines [id:da4809714694919678] 
\draw    (94,268) -- (126,268) ;
%Straight Lines [id:da6937156244672318] 
\draw    (100,264) -- (104,256) ;
%Straight Lines [id:da5268855450380487] 
\draw    (104,264) -- (108,256) ;
%Straight Lines [id:da45382371701169677] 
\draw    (108,264) -- (112,256) ;
%Straight Lines [id:da40135843513989955] 
\draw    (112,264) -- (116,256) ;
%Straight Lines [id:da8797748031397894] 
\draw    (116,264) -- (120,256) ;
%Straight Lines [id:da5396455589833746] 
\draw    (120,264) -- (124,256) ;
%Straight Lines [id:da016938022936106245] 
\draw    (124,264) -- (128,256) ;
%Straight Lines [id:da948682127834509] 
\draw    (98,260) -- (130,260) ;
%Shape: Rectangle [id:dp08846496054139119] 
\draw  [fill={rgb, 255:red, 0; green, 0; blue, 0 }  ,fill opacity=1 ] (100,278) -- (124,278) -- (124,280) -- (100,280) -- cycle ;
%Shape: Rectangle [id:dp5628333633883209] 
\draw  [fill={rgb, 255:red, 0; green, 0; blue, 0 }  ,fill opacity=1 ] (110,272) -- (114,272) -- (114,278) -- (110,278) -- cycle ;
%Straight Lines [id:da04267186502093012] 
\draw    (96,264) -- (128,264) ;

%Shape: Rectangle [id:dp6894530422096654] 
\draw  [fill={rgb, 255:red, 255; green, 255; blue, 255 }  ,fill opacity=1 ][general shadow={fill=black,shadow xshift=0.75pt,shadow yshift=-0.75pt}] (96,208.53) -- (128.15,208.53) -- (128.15,240) -- (96,240) -- cycle ;
%Straight Lines [id:da6601174977334654] 
\draw    (128.15,208) -- (96,241.3) ;

%Straight Lines [id:da8133837980938741] 
\draw [line width=2.25]    (80,192) -- (256,192) ;
%Straight Lines [id:da429414083050879] 
\draw    (112.15,240) -- (112.15,248) ;
%Shape: Rectangle [id:dp9939560108088266] 
\draw  [fill={rgb, 255:red, 255; green, 255; blue, 255 }  ,fill opacity=1 ][general shadow={fill=black,shadow xshift=0.75pt,shadow yshift=-0.75pt}] (152.15,208.53) -- (184.29,208.53) -- (184.29,240) -- (152.15,240) -- cycle ;
%Straight Lines [id:da31690351887272317] 
\draw    (184.29,208) -- (152.15,241.3) ;

%Straight Lines [id:da4925450690668529] 
\draw    (168.15,240) -- (168.15,248) ;
%Straight Lines [id:da9460906042589992] 
\draw    (112,192) -- (112,208) ;
\draw [shift={(112,192)}, rotate = 90] [color={rgb, 255:red, 0; green, 0; blue, 0 }  ][fill={rgb, 255:red, 0; green, 0; blue, 0 }  ][line width=0.75]      (0, 0) circle [x radius= 2.01, y radius= 2.01]   ;
%Straight Lines [id:da7953255510649611] 
\draw    (168,192) -- (168,208) ;
\draw [shift={(168,192)}, rotate = 90] [color={rgb, 255:red, 0; green, 0; blue, 0 }  ][fill={rgb, 255:red, 0; green, 0; blue, 0 }  ][line width=0.75]      (0, 0) circle [x radius= 2.01, y radius= 2.01]   ;
%Straight Lines [id:da20603228224739278] 
\draw    (224,192) -- (224,208) ;
\draw [shift={(224,192)}, rotate = 90] [color={rgb, 255:red, 0; green, 0; blue, 0 }  ][fill={rgb, 255:red, 0; green, 0; blue, 0 }  ][line width=0.75]      (0, 0) circle [x radius= 2.01, y radius= 2.01]   ;
%Shape: Rectangle [id:dp25752390998743535] 
\draw  [fill={rgb, 255:red, 255; green, 255; blue, 255 }  ,fill opacity=1 ][general shadow={fill=black,shadow xshift=0.75pt,shadow yshift=-0.75pt}] (208,207.23) -- (240.15,207.23) -- (240.15,238.7) -- (208,238.7) -- cycle ;
%Straight Lines [id:da8164804894780389] 
\draw    (240.15,206.7) -- (208,240) ;

%Straight Lines [id:da6013771130407253] 
\draw    (224.15,240) -- (224.15,248) ;
%Straight Lines [id:da8978300902586673] 
\draw    (112,168) -- (112,192) ;
%Straight Lines [id:da006449366621711361] 
\draw    (112,112) -- (112,136) ;
%Shape: Triangle [id:dp12868896874136726] 
\draw  [fill={rgb, 255:red, 255; green, 255; blue, 255 }  ,fill opacity=1 ] (112,152) -- (96,168) -- (96,136) -- cycle ;
%Shape: Polygon [id:ds12605022596129922] 
\draw  [fill={rgb, 255:red, 255; green, 255; blue, 255 }  ,fill opacity=1 ][general shadow={fill=black,shadow xshift=0.75pt,shadow yshift=-0.75pt}] (128,136) -- (128,168) -- (96,168) -- (112,152) -- (96,136) -- cycle ;

%Straight Lines [id:da5721175987893949] 
\draw    (112,152) -- (128,152) ;
%Straight Lines [id:da900039405692058] 
\draw    (96,152) -- (64,152) ;

% Text Node
\draw (52.5,129.5) node  [font=\scriptsize] [align=left] {MV-AC};
% Text Node
\draw (105.5,71.5) node  [font=\scriptsize] [align=left] {DC};
% Text Node
\draw (118.5,88.5) node  [font=\scriptsize] [align=left] {AC};
% Text Node
\draw (174.65,88.5) node  [font=\scriptsize] [align=left] {AC};
% Text Node
\draw (161.65,71.5) node  [font=\scriptsize] [align=left] {DC};
% Text Node
\draw (287.5,112.5) node  [font=\scriptsize] [align=left] {AC Microgrid};
% Text Node
\draw (113,8.5) node  [font=\scriptsize] [align=left] {PV};
% Text Node
\draw (168,8.5) node  [font=\scriptsize] [align=left] {BESS};
% Text Node
\draw (224,8.5) node  [font=\scriptsize] [align=left] {AC Loads};
% Text Node
\draw (288.5,192.5) node  [font=\scriptsize] [align=left] {DC Microgrid};
% Text Node
\draw (112,295.5) node  [font=\scriptsize] [align=left] {PV};
% Text Node
\draw (168,295.5) node  [font=\scriptsize] [align=left] {BESS};
% Text Node
\draw (224,295.5) node  [font=\scriptsize] [align=left] {EV Chargers};
% Text Node
\draw (230.65,231.2) node  [font=\scriptsize] [align=left] {DC};
% Text Node
\draw (217.65,214.2) node  [font=\scriptsize] [align=left] {DC};
% Text Node
\draw (174.79,232.5) node  [font=\scriptsize] [align=left] {DC};
% Text Node
\draw (161.79,215.5) node  [font=\scriptsize] [align=left] {DC};
% Text Node
\draw (118.65,232.5) node  [font=\scriptsize] [align=left] {DC};
% Text Node
\draw (105.65,215.5) node  [font=\scriptsize] [align=left] {DC};
% Text Node
\draw (164.5,150.5) node  [font=\scriptsize] [align=left] {Interlinking\\HDT};
% Text Node
\draw (119.5,141.5) node  [font=\scriptsize] [align=left] {AC};
% Text Node
\draw (103,152.5) node  [font=\tiny] [align=left] {AC};
% Text Node
\draw (119.5,161.5) node  [font=\scriptsize] [align=left] {DC};


\end{tikzpicture}}
        \caption{Integration between AC and DC microgrids through a HDT.}
    \end{figure}
\end{frame}

\maketitle